DChain merupakan infrastruktur sertifikat elektronik nasional yang saat ini berbasis \textit{Proof of Authority} (PoA) dan direncanakan bermigrasi ke \textit{Proof of Stake} (PoS) guna meningkatkan desentralisasi. Penelitian ini bertujuan untuk mengevaluasi dan membandingkan kinerja kedua mekanisme konsensus tersebut guna menentukan kelayakan migrasi. Metodologi yang digunakan adalah pendekatan kuantitatif eksperimental melalui \textit{benchmarking} terkomputerisasi menggunakan Hyperledger Caliper pada lingkungan \textit{multi-node} terdistribusi (Docker dan Kubernetes). Pengujian dilakukan melalui berbagai skenario meliputi \textit{throughput}, latensi, efisiensi sumber daya komputasi, dan stabilitas jaringan (\textit{soak test}). Hasil penelitian menunjukkan bahwa \textit{Proof of Stake} (PoS) memiliki kapasitas \textit{throughput} yang lebih tinggi ($\approx$60 TPS) dibandingkan PoA ($\approx$30--40 TPS). Namun, PoA terbukti jauh lebih unggul dalam aspek responsivitas dengan latensi rata-rata 4,6 detik dibandingkan PoS yang mencapai 8 detik, serta memiliki efisiensi penggunaan CPU yang lebih baik. Kesimpulan dari penelitian ini adalah mekanisme konsensus PoA masih menjadi solusi paling optimal untuk kebutuhan DChain saat ini karena karakteristik layanan akademik yang bersifat \textit{user-interactive} dan memerlukan respons cepat, sementara PoS memberikan potensi skalabilitas untuk pengembangan jangka panjang.

\vspace{1cm}

\noindent{Kata kunci} : Blockchain, DChain, Proof of Authority, Proof of Stake, Benchmarking